% Options for packages loaded elsewhere
\PassOptionsToPackage{unicode}{hyperref}
\PassOptionsToPackage{hyphens}{url}
\PassOptionsToPackage{dvipsnames,svgnames,x11names}{xcolor}
%
\documentclass[
  11pt,
  a4paper,
]{ltjsarticle}
\usepackage{amsmath,amssymb}
\usepackage{iftex}
\ifPDFTeX
  \usepackage[T1]{fontenc}
  \usepackage[utf8]{inputenc}
  \usepackage{textcomp} % provide euro and other symbols
\else % if luatex or xetex
  \usepackage{unicode-math} % this also loads fontspec
  \defaultfontfeatures{Scale=MatchLowercase}
  \defaultfontfeatures[\rmfamily]{Ligatures=TeX,Scale=1}
\fi
\usepackage{lmodern}
\ifPDFTeX\else
  % xetex/luatex font selection
\fi
% Use upquote if available, for straight quotes in verbatim environments
\IfFileExists{upquote.sty}{\usepackage{upquote}}{}
\IfFileExists{microtype.sty}{% use microtype if available
  \usepackage[]{microtype}
  \UseMicrotypeSet[protrusion]{basicmath} % disable protrusion for tt fonts
}{}
\makeatletter
\@ifundefined{KOMAClassName}{% if non-KOMA class
  \IfFileExists{parskip.sty}{%
    \usepackage{parskip}
  }{% else
    \setlength{\parindent}{0pt}
    \setlength{\parskip}{6pt plus 2pt minus 1pt}}
}{% if KOMA class
  \KOMAoptions{parskip=half}}
\makeatother
\usepackage{xcolor}
\usepackage[margin=24mm]{geometry}
\usepackage{color}
\usepackage{fancyvrb}
\newcommand{\VerbBar}{|}
\newcommand{\VERB}{\Verb[commandchars=\\\{\}]}
\DefineVerbatimEnvironment{Highlighting}{Verbatim}{commandchars=\\\{\}}
% Add ',fontsize=\small' for more characters per line
\newenvironment{Shaded}{}{}
\newcommand{\AlertTok}[1]{\textcolor[rgb]{1.00,0.00,0.00}{\textbf{#1}}}
\newcommand{\AnnotationTok}[1]{\textcolor[rgb]{0.38,0.63,0.69}{\textbf{\textit{#1}}}}
\newcommand{\AttributeTok}[1]{\textcolor[rgb]{0.49,0.56,0.16}{#1}}
\newcommand{\BaseNTok}[1]{\textcolor[rgb]{0.25,0.63,0.44}{#1}}
\newcommand{\BuiltInTok}[1]{\textcolor[rgb]{0.00,0.50,0.00}{#1}}
\newcommand{\CharTok}[1]{\textcolor[rgb]{0.25,0.44,0.63}{#1}}
\newcommand{\CommentTok}[1]{\textcolor[rgb]{0.38,0.63,0.69}{\textit{#1}}}
\newcommand{\CommentVarTok}[1]{\textcolor[rgb]{0.38,0.63,0.69}{\textbf{\textit{#1}}}}
\newcommand{\ConstantTok}[1]{\textcolor[rgb]{0.53,0.00,0.00}{#1}}
\newcommand{\ControlFlowTok}[1]{\textcolor[rgb]{0.00,0.44,0.13}{\textbf{#1}}}
\newcommand{\DataTypeTok}[1]{\textcolor[rgb]{0.56,0.13,0.00}{#1}}
\newcommand{\DecValTok}[1]{\textcolor[rgb]{0.25,0.63,0.44}{#1}}
\newcommand{\DocumentationTok}[1]{\textcolor[rgb]{0.73,0.13,0.13}{\textit{#1}}}
\newcommand{\ErrorTok}[1]{\textcolor[rgb]{1.00,0.00,0.00}{\textbf{#1}}}
\newcommand{\ExtensionTok}[1]{#1}
\newcommand{\FloatTok}[1]{\textcolor[rgb]{0.25,0.63,0.44}{#1}}
\newcommand{\FunctionTok}[1]{\textcolor[rgb]{0.02,0.16,0.49}{#1}}
\newcommand{\ImportTok}[1]{\textcolor[rgb]{0.00,0.50,0.00}{\textbf{#1}}}
\newcommand{\InformationTok}[1]{\textcolor[rgb]{0.38,0.63,0.69}{\textbf{\textit{#1}}}}
\newcommand{\KeywordTok}[1]{\textcolor[rgb]{0.00,0.44,0.13}{\textbf{#1}}}
\newcommand{\NormalTok}[1]{#1}
\newcommand{\OperatorTok}[1]{\textcolor[rgb]{0.40,0.40,0.40}{#1}}
\newcommand{\OtherTok}[1]{\textcolor[rgb]{0.00,0.44,0.13}{#1}}
\newcommand{\PreprocessorTok}[1]{\textcolor[rgb]{0.74,0.48,0.00}{#1}}
\newcommand{\RegionMarkerTok}[1]{#1}
\newcommand{\SpecialCharTok}[1]{\textcolor[rgb]{0.25,0.44,0.63}{#1}}
\newcommand{\SpecialStringTok}[1]{\textcolor[rgb]{0.73,0.40,0.53}{#1}}
\newcommand{\StringTok}[1]{\textcolor[rgb]{0.25,0.44,0.63}{#1}}
\newcommand{\VariableTok}[1]{\textcolor[rgb]{0.10,0.09,0.49}{#1}}
\newcommand{\VerbatimStringTok}[1]{\textcolor[rgb]{0.25,0.44,0.63}{#1}}
\newcommand{\WarningTok}[1]{\textcolor[rgb]{0.38,0.63,0.69}{\textbf{\textit{#1}}}}
\setlength{\emergencystretch}{3em} % prevent overfull lines
\providecommand{\tightlist}{%
  \setlength{\itemsep}{0pt}\setlength{\parskip}{0pt}}
\setcounter{secnumdepth}{5}
\ifLuaTeX
  \usepackage{selnolig}  % disable illegal ligatures
\fi
\IfFileExists{bookmark.sty}{\usepackage{bookmark}}{\usepackage{hyperref}}
\IfFileExists{xurl.sty}{\usepackage{xurl}}{} % add URL line breaks if available
\urlstyle{same}
\hypersetup{
  colorlinks=true,
  linkcolor={blue},
  filecolor={Maroon},
  citecolor={Blue},
  urlcolor={blue},
  pdfcreator={LaTeX via pandoc}}

\author{}
\date{}

\begin{document}

\hypertarget{ux7ddaux5f62ux4ee3ux6570ux306eux5168ux4f53ux5730ux56f3}{%
\section{線形代数の全体地図}\label{ux7ddaux5f62ux4ee3ux6570ux306eux5168ux4f53ux5730ux56f3}}

線形代数を一言で表すと、\textbf{ベクトル空間と、その間を結ぶ線形写像を研究する学問}です。行列はその線形写像を座標で表現したものにすぎません。

\hypertarget{ux3064ux306eux4e2dux5fc3ux30c6ux30fcux30de}{%
\subsection{3つの中心テーマ}\label{ux3064ux306eux4e2dux5fc3ux30c6ux30fcux30de}}

\hypertarget{ux7a7aux9593ux306eux69cbux9020}{%
\subsubsection{1. 空間の構造}\label{ux7a7aux9593ux306eux69cbux9020}}

ベクトル、一次結合、span、一次独立、基底、次元を通して、「どの方向が本当に独立な情報なのか」を調べます。

\hypertarget{ux7ddaux5f62ux5199ux50cfux306eux69cbux9020}{%
\subsubsection{2.
線形写像の構造}\label{ux7ddaux5f62ux5199ux50cfux306eux69cbux9020}}

\texttt{Ax}
は単なる行列計算ではなく、ある空間から別の空間への写像です。核
\texttt{ker(A)} は消えてしまう方向、像 \texttt{Im(A)}
は実際に到達できる方向です。

\hypertarget{ux826fux3044ux5ea7ux6a19ux7cfbux3067ux898bux308b}{%
\subsubsection{3.
良い座標系で見る}\label{ux826fux3044ux5ea7ux6a19ux7cfbux3067ux898bux308b}}

同じ線形写像でも、基底を変えると行列の見え方が変わります。固有値分解や
SVD は、線形変換を理解しやすい座標系へ移してから観察する方法です。

\hypertarget{ux5b66ux7fd2ux306eux5927ux304dux306aux6d41ux308c}{%
\subsection{学習の大きな流れ}\label{ux5b66ux7fd2ux306eux5927ux304dux306aux6d41ux308c}}

\begin{Shaded}
\begin{Highlighting}[]
\NormalTok{ベクトル空間}
\NormalTok{  ↓}
\NormalTok{基底・次元}
\NormalTok{  ↓}
\NormalTok{線形写像・行列}
\NormalTok{  ↓}
\NormalTok{核・像・ランク}
\NormalTok{  ↓}
\NormalTok{内積・直交・射影}
\NormalTok{  ↓}
\NormalTok{固有値・固有ベクトル}
\NormalTok{  ↓}
\NormalTok{行列分解}
\NormalTok{  ├─ LU}
\NormalTok{  ├─ QR}
\NormalTok{  ├─ 固有値分解}
\NormalTok{  └─ SVD}
\NormalTok{       ↓}
\NormalTok{   最小二乗・擬似逆}
\NormalTok{       ↓}
\NormalTok{   PCA・低ランク近似}
\end{Highlighting}
\end{Shaded}

\hypertarget{svd-ux306fux3069ux3053ux306bux4f4dux7f6eux3059ux308bux304b}{%
\subsection{SVD
はどこに位置するか}\label{svd-ux306fux3069ux3053ux306bux4f4dux7f6eux3059ux308bux304b}}

SVD

\[
A=U\Sigma V^\top
\]

は、入力側の直交基底 \texttt{V}、各方向の伸縮量
\texttt{Σ}、出力側の直交基底 \texttt{U} に線形写像を分解します。

したがって SVD を本当に理解するには、少なくとも次が必要です。

\begin{itemize}
\tightlist
\item
  基底と座標
\item
  線形写像
\item
  ランク
\item
  内積と直交
\item
  固有値と対称行列
\end{itemize}

SVD
は線形代数の入口ではなく、前半で学んだ概念が一つに集約される地点だと考えてください。

\hypertarget{ux5230ux9054ux76eeux6a19}{%
\subsection{到達目標}\label{ux5230ux9054ux76eeux6a19}}

最終的には、公式を見て計算するだけでなく、行列を見たときに

\begin{itemize}
\tightlist
\item
  何次元の情報が保存されるか
\item
  どの方向が消えるか
\item
  どの方向が大きく伸びるか
\item
  逆問題が安定か
\item
  低次元近似が可能か
\end{itemize}

を幾何学的に説明できることを目標にします。

\end{document}
